\section{Introduction}

\subsection{The Main Results}

We consider the incompressible Euler equations in the plane, in vorticity form:
\begin{equation}
\label{eq:euler}
\partial_t \omega + v\cdot\nabla \omega = 0,\qquad
v = \KBS\ast \omega,\qquad
\omega(t=0) = \omega_0 .
\end{equation}
Here $\KBS(x)=x^\perp/(2\pi|x|^2)$, $x^\perp=(-x_2,x_1)$, denotes the
Biot-Savart kernel, so that $v$ is divergence free and
$\partial_1v_2-\partial_2v_1=\omega$. There is no external force, and we call
$\omega_0$ the datum.

This paper is the second of two. Together with Part~I~\cite{PartI}, it proves
that, without force in the plane, uniqueness fails among renormalized weak
solutions with vorticity in $L^\infty_t(L^1\cap L^q)_x$, for every
$q\in(4/3,\infty)$ at once and for an explicit radial datum. Numbers with the
prefix I refer to statements, sections and equations of Part~I. Section~1 of
Part~I states the question and discusses the literature, including the
non-uniqueness with a radial force of Vishik \cite{Vishik1}, \cite{Vishik2}, in
the account of \cite{ABCDGJK}, and the results obtained without force by convex
integration. It also describes the construction heuristically. We do not
repeat this discussion.

For $T>0$ and $1\le q\le\infty$ we write $X_q(T)$ for the space
$L^\infty_t(L^1\cap L^q)_x$ of functions on $(0,T)\times\R^2$. Renormalized weak solutions in $X_q(T)$, $4/3<q<\infty$ (after
DiPerna and Lions \cite{DiPernaLions}), and equimeasurability are defined in
Definitions~I.\ref{I-def:weak} and~I.\ref{I-def:equimeasurable}. The ring $\Phi$ of
Definition~I.\ref{I-def:ring} is a closed formula in $\log|y|$ times a fixed smooth
cutoff; it is radial, $\Phi\in C_c^\infty(\R^2)$, $0\le\Phi<1$,
$\Phi\not\equiv0$, and $\supp\Phi\subset\{a\le|y|\le1\}$ for an explicit
$a\in(0,1)$ (Lemma~I.\ref{I-lem:ring-facts}). For $A_1>0$, $\mu\in(1,2]$ and
$\sbar\in(0,a/16]$ we put
\begin{equation}
\label{eq:datum}
\begin{gathered}
f=\sum_{k\ge1}A_k\,\Phi(\cdot/r_k),\qquad A_k=A_1\mu^{k-1},\\
r_1=1,\qquad r_{k+1}=s_kr_k,\qquad s_k=\sbar(2\mu)^{-k}\qquad(k\ge1).
\end{gathered}
\end{equation}
This is Definition~\ref{def:datum}. We call $A_k\Phi(\cdot/r_k)$ ring $k$ of
$f$. It is supported in the annulus $\{ar_k\le|x|\le r_k\}$, and these annuli
are pairwise disjoint (Proposition~\ref{prop:integrability}). The main theorem
of the two papers is the following; it is Theorem~I.\ref{I-thm:main}, restated.

\begin{theorem}
\label{thm:main}
Let $\Phi$ and $a$ be as in Definition~I.\ref{I-def:ring}. There are
$\mu\in(1,\frac{17}{16}]$ and $\sbar\in(0,a/16]$ such that, for every $A_1>0$,
the function $f$ of Equation~\eqref{eq:datum} has the following properties. It
is radial and nonnegative, $f\in (L^1\cap L^q)(\R^2)$ for every $q\in[1,\infty)$,
and $f\notin L^\infty(\R^2)$. There is $\Tstar\in(0,\infty)$ such that, for
every $T>\Tstar$ and every $q\in(4/3,\infty)$, Equation~\eqref{eq:euler} with
$\omega_0=f$ has two distinct renormalized weak solutions in $X_q(T)$, in the
sense of Definition~I.\ref{I-def:weak}:
\begin{enumerate}[label=\textup{(\roman*)}]
\item the stationary solution $\omega(t)=f$, $t\in[0,T]$;
\item a solution $\omega_\infty$ such that
  $\omega_\infty\in C([0,\infty);L^p(\R^2))$ for every $p\in[1,\infty)$ and,
  for every $t\ge0$, $\omega_\infty(t)$ is equimeasurable with $f$ and
  $3$-fold rotationally symmetric.
\end{enumerate}
More precisely, there is $t_0\in[0,\Tstar)$ such that $\omega_\infty$ and $f$
are distinct elements of $X_q(T')$ for every $T'>t_0$ and every
$q\in(4/3,\infty)$.
\end{theorem}

We do not treat $q\le4/3$. No symmetry is
imposed on the class; the $3$-fold rotational symmetry of $\omega_\infty$ comes from the
construction, where it removes the angular modes $\ell\notin3\Z$, among them the
modes $\ell=\pm1$ that contain the translations of a ring, and leaves the
modes $\ell\in3\Z$, whose spectra Part~I controls under
Hypothesis~I.\ref{I-hyp:N}.

The proof of Theorem~\ref{thm:main} is divided between Part~I and this paper.
Part~I proves a non-uniqueness criterion (Section~I.\ref{I-sec:compact}) and the
linear analysis of the ring (Section~I.\ref{I-sec:ring}), and it uses no machine
computation. The main conclusions of its linear analysis assume three explicit
statements, collected in Section~I.\ref{I-sec:ring} as
Hypothesis~I.\ref{I-hyp:N}. The statements \textup{(N1)} and \textup{(N2)} locate
the zeros in the mode $\ell=3$ of the Evans function of the ring, an explicit
holomorphic function, for every $G\in[0,\frac1{50}]$, where $2\pi G$ is the
circulation of a point vortex at the origin (in the reference units of
Definition~I.\ref{I-def:ring}; this is $G\in[0,\Gone]$ in ring units). The
statement \textup{(N3)} is an inequality between explicit numbers controlling
the modes $\ell=\pm6$ and $|\ell|\ge9$. This paper proves that
\textup{(N1)}--\textup{(N3)} imply Theorem~\ref{thm:main}, by a nonlinear
construction followed by the criterion of Part~I, and it proves
\textup{(N1)}--\textup{(N3)}. The implication is the following reduction,
which we prove in Section~\ref{sec:bootstrap}.

\begin{proposition}
\label{prop:reduction}
Suppose that \textup{(N1)}--\textup{(N3)} of Hypothesis~I.\ref{I-hyp:N} hold. Then
Theorem~\ref{thm:main} holds.
\end{proposition}

The proof of Proposition~\ref{prop:reduction} uses no machine computation.
Appendix~\ref{app:verify} proves \textup{(N3)} by hand and verifies
\textup{(N1)} and \textup{(N2)} with rigorous interval arithmetic, the only
machine-computed part of the proof (code: the package
\texttt{euler-ring-certificate}, version~1.0.0, in the public repository
\url{https://github.com/mcoiculescu/NonUniqueEuler};
Section~\ref{app:reproduce}). It thus proves the following theorem.

\begin{theorem}
\label{thm:N-holds}
Hypothesis~I.\ref{I-hyp:N} holds: the statements \textup{(N1)}, \textup{(N2)} and
\textup{(N3)} are true.
\end{theorem}

\begin{proof}[Proof of Theorem~\ref{thm:main}]
By Theorem~\ref{thm:N-holds}, the hypotheses \textup{(N1)}--\textup{(N3)}
hold, so Theorem~\ref{thm:main} follows from Proposition~\ref{prop:reduction}.
\end{proof}

\subsection{Results of Part I}
\label{sec:II-partI}

Part~I's introduction summarizes what this paper uses in two theorems: the
conditional non-uniqueness theorem (Theorem~I.\ref{I-thm:conditional}, the form
of Corollary~I.\ref{I-cor:defect} given by Remark~I.\ref{I-rem:conditional}) and
the linear instability of the ring under \textup{(N1)}--\textup{(N3)}
(Theorem~I.\ref{I-thm:linear}).

In detail, we use the following definitions and results of Part~I, which we
cite by their numbers there and do not restate. From
Section~I.\ref{I-sec:prelim}, we use the properties of the Biot-Savart kernel
(Lemma~I.\ref{I-lem:BS-basic}) and the multipoles, equimeasurability and
renormalized weak solutions (Definitions~I.\ref{I-def:multipole},
I.\ref{I-def:equimeasurable} and~I.\ref{I-def:weak}, with
Equations~(I.\ref{I-eq:weak}) and~(I.\ref{I-eq:renorm})). From
Section~I.\ref{I-sec:compact}, we use the fact that a radial function in
$(L^1\cap L^q)(\R^2)$, $4/3<q<\infty$, is a stationary renormalized weak
solution (Lemma~I.\ref{I-lem:radial-stationary}), the $2$-admissible sequences
(Definition~I.\ref{I-def:adm}), and the non-uniqueness criterion of
Corollary~I.\ref{I-cor:defect}, which rests on Theorem~I.\ref{I-thm:criterion}
and which we apply in the proof of Proposition~\ref{prop:reduction}. From
Section~I.\ref{I-sec:ring}, we use the ring and the annuli $N$, $N'$
(Definition~I.\ref{I-def:ring}, Lemma~I.\ref{I-lem:ring-facts}); the mode
settings, the decomposition into angular modes, the Evans function and its
simple zeros, and the semigroup (Definition~I.\ref{I-def:sector-setting},
Lemmas~I.\ref{I-lem:sector}, I.\ref{I-lem:evans-basic}
and~I.\ref{I-lem:semigroup-construct}, Theorem~I.\ref{I-thm:evans-simple}); and
Hypothesis~I.\ref{I-hyp:N}. Under \textup{(N1)}--\textup{(N3)} we use, from the
same section, the spectrum of the linearization (Proposition~I.\ref{I-prop:SU}),
the eigendata (Definition~I.\ref{I-def:eigendata}), the continuity in $G$ of
$\lamst(G)$ and of the sizes $|\pout_*(G)|$, $|\beta(G)|$,
$\|\ellst(G)\|_{(X^0)^*}$ and $\|\phist(G)\|_{H^k}$ of the eigendata, the
semigroup bounds and the uniformity of the constants in $G$
(Propositions~I.\ref{I-prop:U} and~I.\ref{I-prop:semigroup},
Lemma~I.\ref{I-lem:uniform}).

\subsection{Strategy}

We describe the construction with references to the statements of this paper
that carry it out. The background is the datum $f$, which is radial, hence
stationary (Lemma~I.\ref{I-lem:radial-stationary}). Its inner rings have larger
amplitudes and rotate faster, and the task is to
carry an instability from an inner ring outward through every ring in a time
independent of the depth, uniformly in the number of interacting rings.

The linear input is the instability of one ring. The linearization of
Equation~\eqref{eq:euler} at $\Phi$ plus a point vortex of small circulation
$2\pi G$ at the origin, which stands for everything inside the ring, is the
operator $\Lin(G)$ of Theorem~I.\ref{I-thm:linear}. On $3$-fold rotationally
symmetric perturbations it has an eigenvalue $\lamst(G)$ of the mode $\ell=3$,
algebraically simple there, with real part $\sigma(G)>0$, and its conjugate
$\overline{\lamst(G)}$ in the mode $\ell=-3$; the rest of the spectrum has
smaller real part (Proposition~I.\ref{I-prop:SU}). Ring $k$ of $f$ evolves on
the time scale $1/A_k$ and feels the inner rings mostly through their
circulation $2\pi G_kA_kr_k^2$, with $G_k\ge0$ small when $\sbar$ is small
(Lemma~\ref{lem:inner-circ}). In ring-$k$ units (Section~\ref{sec:datum}), the linearization at
ring $k$ plus this point vortex is therefore $\Lin(G_k)$.

The approximate datum $f^{(n)}$ keeps the first $n$ rings of $f$ and replaces
ring $n$, the innermost of them, by a rearrangement of it whose coefficient
along the unstable eigenfunction has a small size of order $\kseed$, fixed
independently of $n$ (Definition~\ref{def:approx}, Lemma~\ref{lem:seed}). As
$r_n$ decreases faster than $A_n$ increases, $f^{(n)}\to f$ in
$(L^1\cap L^q)(\R^2)$ for every $q<\infty$ (Lemma~\ref{lem:adm}). Each
$f^{(n)}$ is smooth with compact support, so it has a unique global smooth
solution $\omega^{(n)}$ (Theorem~\ref{thm:smooth}), which is $3$-fold
rotationally symmetric (Lemma~\ref{lem:adm}).

In $\omega^{(n)}$ the instability passes outward like a relay. The unstable
coefficient of ring $k$ grows at the rate $A_k\sigma(G_k)$ to a fixed size
$\delta_0$, and we then say that ring $k$ saturates. Its mode-$3$ moment drives
ring $k-1$ through a field smaller by a factor of order $s_{k-1}^5$
(Lemma~\ref{lem:coupling}). This drive outgrows the free growth of ring $k-1$
by a margin of order $\sigma(0)(\mu-1)/\mu$. The forced part of the
coefficient of ring $k-1$ is a truncated Laplace transform of the drive
(Lemma~\ref{lem:coupling}, Equation~\eqref{eq:kick-A}). Because of this
margin, it is dominated by the end of the growth and pinned after saturation. Its modulus lies between two multiples of the size of
the drive at saturation divided by $\mu-1$, with constants independent of $k$
and $n$ (Proposition~\ref{prop:kick} and its proof).

As $\mu$ is fixed, ring $k-1$ therefore starts with a coefficient of order
$s_{k-1}^5$ (Step~3 in the proof of Proposition~\ref{prop:bootstrap}), and it
saturates after a physical time of about
$5\log(1/s_{k-1})/(\sigma(G_{k-1})A_{k-1})$. Since $\log(1/s_k)$ grows
linearly in $k$ and $A_k$ geometrically, the total time is bounded by a
$\Tstar$ independent of $n$ (Lemma~\ref{lem:time}).

The phase clock of Definition~\ref{def:clock} fixes, from $\omega^{(n)}$
itself, the time at which each ring becomes active and the time at which it is
expected to saturate. Until it becomes active, a ring stays close to $\Phi$,
although forced by the fields inside and outside it (Theorem~\ref{thm:quiet}).
The active ring follows its unstable eigenfunction (Theorem~\ref{thm:active}).
Once the ring just inside it has saturated, the material inside each ring
stays in a small disk (Proposition~\ref{prop:confine}). The hypotheses of these three statements are
verified jointly, for all rings and all times until the outermost ring
saturates, by one bootstrap with constants independent of $n$ (proof of
Proposition~\ref{prop:bootstrap}).

When the outermost ring saturates, at $T^{(n)}\le\Tstar$, its moment
$\int\omega^{(n)}(T^{(n)})\Psione$, where $\Psione(x)=\chiout(|x|)\,\bar x^{\,3}$
with a fixed real cutoff $\chiout\in C_c^\infty((a/4,3))$, has modulus at least
$c_0>0$ independent of $n$ (proof of Proposition~\ref{prop:reduction}), while
it vanishes for $f$. This lower bound is the moment hypothesis of
Theorem~I.\ref{I-thm:conditional}. The $\omega^{(n)}$ are bounded in
$(L^1\cap L^2)(\R^2)$ uniformly in time, and compactness of renormalized
solutions in $L^2$ (Section~I.\ref{I-sec:compact}) gives a subsequence
converging to a renormalized weak solution $\omega_\infty$ with datum $f$. The
moment bound passes to the limit and separates $\omega_\infty$ from $f$
(Corollary~I.\ref{I-cor:defect}).

Section~\ref{sec:II-prelim} collects notation and preliminaries.
Section~\ref{sec:single} studies, in ring units, one ring of the approximate
solutions, driven by the field of the vorticity inside and outside it: the quiet
ring, the active ring and the seeded ring.
Section~\ref{sec:datum} defines the datum, the approximants $f^{(n)}$ and the
phase clock. Section~\ref{sec:several} bounds the fields that the rings induce
on one another, and it confines the material inside a ring after the ring just
inside it has saturated. Section~\ref{sec:kick} computes the kick that a ring
receives from the rings inside it. Section~\ref{sec:relay} proves the relay
lemmas, which are the estimates on the rings that the bootstrap combines, each
under explicit hypotheses; it also proves the time bound of Lemma~\ref{lem:time}.
Section~\ref{sec:param} chooses the parameters so that the conditions of the
relay lemmas on them hold; the choices are made in the order of
Table~\ref{tab:param}, and none of them depends on $n$, $A_1$ or
$\omega^{(n)}$ (Lemma~\ref{lem:nonempty}). Section~\ref{sec:bootstrap} proves
the joint bootstrap (Proposition~\ref{prop:bootstrap}) and
Proposition~\ref{prop:reduction}. Appendices~\ref{app:verify}
and~\ref{app:smooth} prove Theorems~\ref{thm:N-holds} and~\ref{thm:smooth}, respectively.
Appendix~\ref{app:deps} lists, for each statement on the path to
Theorem~\ref{thm:main}, the statements of this paper and of Part~I that its
proof uses and the parts of \textup{(N1)}--\textup{(N3)} on which it depends.

The main idea of a linear instability in physical coordinates driving the
non-uniqueness mechanism, the use of non-monotone unstable radial vortices and of
a multiscale datum, and the overall strategy of the proof are due to the author.
Within this strategy, the detailed proofs and estimates, the code for the
computer-assisted verification and the text were largely produced with
Claude Opus 5.5 (Anthropic), under the direction of the author. The author has
checked all proofs by hand and takes full responsibility for the paper.

\paragraph{Acknowledgments.} I am supported at New York University by a Junior
Fellowship of the Simons Society of Fellows, funded by Simons Foundation
International and administered by the Simons Foundation. I would also like to thank
Tristan Buckmaster for generously offering his time and advice on this and many
other projects.

\subsection{Preliminaries}
\label{sec:II-prelim}

This subsection collects the notation that we take from Part~I and the
preliminaries used only in this paper. We use all notation and conventions of
Section~I.\ref{I-sec:prelim}, among them the convention on constants and
$O(\cdot)$, which we follow. From
Section~I.\ref{I-sec:ring} we use the convention $m=3$, the ring units and the
number $\Gone$ of Definition~I.\ref{I-def:ring}, the annuli $N$ and $N'$ of
Section~I.\ref{I-sec:ring-def}, the spaces $X^s$ and the operators $\Lin(G)$ on
them, fixed before Proposition~I.\ref{I-prop:SU}, and the constants $c_N$, $C_N$,
$\cout$, $\|\ellst\|$, $\Cls$, $\Cp$ and $\Cnab(k)$ of
Lemma~I.\ref{I-lem:uniform}. The results and definitions of
Part~I that we use are listed in Section~\ref{sec:II-partI}.

In ring units we also use the closed annuli
\[
\NNm=\{\tfrac78a\le|y|\le\tfrac54\},\qquad \NN=\{\tfrac34a\le|y|\le\tfrac32\},\qquad
\Nflat=\{\tfrac58a\le|y|\le\tfrac74\}.
\]
They mark where the support of a ring lies at the stages of the construction:
in $\NNm$ for a quiet ring and at the start of the active phase
(Theorems~\ref{thm:quiet} and~\ref{thm:active}), in $\NN$ during the active
phase, and in $\Nflat$ under the bootstrap assumption of
Section~\ref{sec:bootstrap}, which the bootstrap improves to $\NN$. All of them
lie in $N$, the common support of the spaces $X^s$ of Part~I, and the strict
inclusion $\NN\subset\operatorname{int}\Nflat$ leaves the room that the bootstrap
needs.

\begin{lemma}
\label{lem:annuli}
Let $a\in(0,1)$ and define the annuli above, and $N$, $N'$, with this $a$. Then
$\{a\le|y|\le1\}\subset\NNm\subset\NN\subset\operatorname{int}\Nflat\subset\Nflat\subset N\subset N'$.
Let $S\subset\{a\le|y|\le1\}$ \textup{(}respectively $S\subset\NNm$\textup{)},
and let $S'$ be a set such that every $y'\in S'$ satisfies
$\big||y'|-|y|\big|\le a/16$ for some $y\in S$. Then $S'\subset\NNm$
\textup{(}respectively $S'\subset\NN$\textup{)}.
\end{lemma}
\begin{proof}
The inner radii $a$, $\frac78a$, $\frac34a$, $\frac58a$, $\frac12a$, $\frac14a$
strictly decrease and the outer radii $1$, $\frac54$, $\frac32$, $\frac74$, $2$, $3$
strictly increase, and the last two radii of each list are those of $N$ and $N'$;
this gives the chain. If $a\le|y|\le1$ and
$||y'|-|y||\le a/16$, then $\frac{15}{16}a\le|y'|\le1+\frac a{16}$, and
$\frac{15}{16}a\ge\frac78a$ and $1+\frac a{16}\le\frac54$ because $a\le4$. If
instead $y\in\NNm$, then $\frac{13}{16}a\le|y'|\le\frac54+\frac a{16}$, and
$\frac{13}{16}a\ge\frac34a$ and $\frac54+\frac a{16}\le\frac32$, for the same
reason.
\end{proof}

Away from the support of a vorticity, complex variables turn the Biot-Savart
law into a Cauchy integral, which can be expanded in powers of the complex
variable.
We call $V_1-iV_2$ the \textbf{complex form} of a vector field $V$, and wherever
we use complex forms we write $z=x_1+ix_2$ (this $z$ is not the number
$\lamst(0)$ of Definition~I.\ref{I-def:eigendata}). Since $(\KBS)_1(x)-i(\KBS)_2(x)=(-x_2-ix_1)/(2\pi|x|^2)=1/(2\pi iz)$,
the velocity of every $g\in L^1(\R^2)$ has the complex form
\begin{equation}
\label{eq:BS-complex}
(\KBS\ast g)_1(z)-i(\KBS\ast g)_2(z)=\frac1{2\pi i}\int\frac{g(\zeta)}{z-\zeta}\dd A(\zeta)
\qquad(z\in\C\setminus\supp g),
\end{equation}
where $\dd A$ is the area element, the integral converges absolutely and, by differentiation under the
integral sign, defines a holomorphic function of $z\in\C\setminus\supp g$. Let
$V$ be a real vector field on an open set $U$ whose complex form $F$ is
holomorphic on $U$. Then $\partial_{x_1}F=F'$ and $\partial_{x_2}F=iF'$, so every
partial derivative of order $q$ of $F$ is $i^{\gamma_2}F^{(q)}$ for the
corresponding multi-index $\gamma$. Fix one of the $2^q$ ordered $q$-tuples of
directions. Since $V_1=\Re F$ and $V_2=-\Im F$, the corresponding derivatives of
$V_1$ and $V_2$ at $x$ are the real part and minus the imaginary part of the
corresponding derivative of $F$ at $z$, a number of modulus $|F^{(q)}(z)|$; so this
pair of entries of $\nabla^qV(x)$ has Euclidean norm $|F^{(q)}(z)|$. Summing the
squares over the $2^q$ tuples gives
\begin{equation}
\label{eq:BS-complex-deriv}
|\nabla^qV(x)|=2^{q/2}\,|F^{(q)}(z)|\qquad(z\in U,\ q\ge0).
\end{equation}

The next two lemmas bound the field of a rotationally symmetric vorticity away
from its support: in a disk inside a hole of the support, and outside a disk that
contains the support, after the monopole is removed. Section~\ref{sec:several}
applies them to the rings outside and inside a given ring, respectively.

\begin{lemma}
\label{lem:hollow}
Let $m_0\ge2$ be an integer and $R>0$. Let $\nu\in (L^1\cap L^\infty)(\R^2)$ be
real and $m_0$-fold rotationally symmetric with $\supp\nu\subset\{|y|\ge R\}$, and let
$\bar\nu(y)=\frac1{2\pi}\int_0^{2\pi}\nu(\Orot_\phi y)\dd\phi$ be its angular
average. Then $\KBS\ast\bar\nu=0$ on $B_R$, and for $|x|\le R/2$,
\[
|\KBS\ast\nu(x)|\le\tfrac4{3\pi}\|\nu\|_{L^1}|x|^{m_0-1}R^{-m_0}.
\]
\end{lemma}

\begin{proof}
Let $|z|<R$, so that $|z-\zeta|\ge R-|z|>0$ for $\zeta\in\supp\nu$, and let $u=\KBS\ast\nu$. For $|\zeta|\ge R$ the series
$(z-\zeta)^{-1}=-\sum_{l\ge1}z^{l-1}\zeta^{-l}$ converges uniformly in $\zeta$, with ratio
at most $|z|/R$, so Equation~\eqref{eq:BS-complex} gives, term by term,
$u_1-iu_2=\frac{i}{2\pi}\sum_{l\ge1}c_l(\nu)z^{l-1}$ with
$c_l(\nu)=\int\nu(\zeta)\zeta^{-l}\dd A(\zeta)$. The substitution $\zeta\mapsto e^{2\pi i/m_0}\zeta$ and the $m_0$-fold rotational
symmetry of $\nu$ show $c_l(\nu)=e^{-2\pi
il/m_0}c_l(\nu)$, so $c_l(\nu)=0$ unless $m_0\mid l$. The angular average
$\bar\nu$ is radial and, like $\nu$, supported in $\{|y|\ge R\}$, so the same
expansion applies to it, and integration in the angle
shows $c_l(\bar\nu)=0$ for $l\ge1$; this proves $\KBS\ast\bar\nu=0$ on $B_R$. It
also gives $c_l(\nu)=c_l(\nu-\bar\nu)$, so $|c_l(\nu)|\le\|\nu-\bar\nu\|_{L^1}R^{-l}$,
as $|\zeta|\ge R$ on the support of $\nu-\bar\nu$.
Since $\nu$ is real, $|u|=|u_1-iu_2|$. Only the terms with $l=jm_0$, $j\ge1$,
remain, so for $|x|\le R/2$ these bounds and the sum of a geometric series
give
\[
|u(x)|\le\frac{\|\nu-\bar\nu\|_{L^1}}{2\pi}\sum_{j\ge1}|x|^{jm_0-1}R^{-jm_0}
\le\frac{\|\nu-\bar\nu\|_{L^1}}{2\pi(1-2^{-m_0})}|x|^{m_0-1}R^{-m_0},
\]
and $1/(1-2^{-m_0})\le\frac43$ because $m_0\ge2$. Finally
$\|\nu-\bar\nu\|_{L^1}\le2\|\nu\|_{L^1}$, because averaging over rotations does
not increase the $L^1$ norm.
\end{proof}

For an integer $m\ge1$ and $g$ as in Definition~I.\ref{I-def:multipole}, let
$\pout_m[g]$ and $\Qm[g]$ be the multipoles defined there.
If $\supp g\subset\overline B_{\varrho}$, the expansion
\[
\log|x-y|=\log|x|-\Re\sum_{k\ge1}\frac1k\Big(\frac yx\Big)^k\quad(|y|<|x|)
\]
shows that the angular mode $m$ of the stream function
$(2\pi)^{-1}\log|\cdot|\ast g$ of $g$ equals $-(2m)^{-1}\pout_m[g]\,r^{-m}$
for $r>\varrho$; the series may be integrated term by term against $g$ because,
for fixed $|x|>\varrho$, it converges uniformly in $y\in\overline B_\varrho$. For real $g$ this is the angular mode $m$ of the exterior
expansion after Definition~I.\ref{I-def:multipole}.
Section~\ref{sec:kick} uses it for the field that the inner rings induce on a
ring.

\begin{lemma}
\label{lem:ext-multipole}
Let $m\ge1$ and $q\ge0$ be integers, and let $\varrho>0$. Let $g\in L^1(\R^2)$ be
real and $m$-fold rotationally symmetric with $\supp g\subset\overline{B}_\varrho$, and let $E$ be the
velocity field $\KBS\ast g$ minus its monopole
$\frac{1}{2\pi}\big(\int g\big)\frac{x^\perp}{|x|^2}$. Then for $|x|\ge2\varrho$,
\[
|\nabla^qE(x)|\le C_{m,q}\,\|g\|_{L^1}\,\varrho^m|x|^{-m-1-q},\qquad
C_{m,q}=\frac{2^q}{2\pi}\sup_{0<\rho\le1/2}\rho^{-m}\sum_{l\in m\N}(l+q)^q\rho^{l}<\infty .
\]
\end{lemma}

\begin{proof}
Let $|z|>\varrho$. For $|\zeta|\le\varrho$ the series
$(z-\zeta)^{-1}=\sum_{l\ge0}\zeta^lz^{-l-1}$ converges uniformly in $\zeta$, so
Equation~\eqref{eq:BS-complex} gives
$(\KBS\ast g)_1-i(\KBS\ast g)_2=\frac1{2\pi i}\sum_{l\ge0}Q_lz^{-l-1}$ with
$Q_l=\int g(\zeta)\zeta^l\dd A(\zeta)$. The term $l=0$, which is $(\int g)/(2\pi iz)$, is the
monopole, and $Q_l=0$ unless $m\mid l$, because the substitution $\zeta\mapsto
e^{2\pi i/m}\zeta$ and the $m$-fold rotational symmetry of $g$ give $Q_l=e^{2\pi il/m}Q_l$. Hence the complex form of $E$ is
$F(z)=\frac1{2\pi i}\sum_{l\in m\N}Q_lz^{-l-1}$, where
$|Q_l|\le\|g\|_{L^1}\varrho^l$ because $|\zeta|\le\varrho$ on $\supp g$. The series converges locally uniformly on
$\{|z|>\varrho\}$, so $F$ is holomorphic there, may be differentiated term by
term, and, as $E$ is real,
Equation~\eqref{eq:BS-complex-deriv} gives
$|\nabla^qE(x)|=2^{q/2}|F^{(q)}(z)|\le2^q|F^{(q)}(z)|$. With
$|\partial_z^qz^{-l-1}|=(l+1)\cdots(l+q)|z|^{-l-1-q}\le(l+q)^q|z|^{-l-1-q}$ and
$\rho=\varrho/|z|\le\frac12$,
\[
|F^{(q)}(z)|\le\frac{\|g\|_{L^1}}{2\pi}|z|^{-1-q}\sum_{l\in m\N}(l+q)^q\rho^l
\le\frac{C_{m,q}}{2^q}\|g\|_{L^1}\rho^m|z|^{-1-q},
\]
and $\rho^m|z|^{-1-q}=\varrho^m|x|^{-m-1-q}$. The constant is finite, because
$\rho^{-m}\sum_{l\in m\N}(l+q)^q\rho^l=\sum_{j\ge1}(mj+q)^q\rho^{m(j-1)}$ is
nondecreasing in $\rho$ and finite at $\rho=\frac12$.
\end{proof}

The approximate solutions $\omega^{(n)}$ of Definition~\ref{def:approx} are the
global smooth solutions with the smooth, compactly supported data $f^{(n)}$,
given by the following classical theorem, which is proved in
Appendix~\ref{app:smooth}. Besides the properties (1)--(3) of the solution, the
proof of Theorem~\ref{thm:main} uses the flow of item~(4), the uniqueness of
item~(5), which identifies a rescaled or time-translated solution with the
solution of its datum, the symmetry of item~(6), and item~(7), by which these solutions are renormalized weak solutions
in the sense of Part~I.

\begin{theorem}[cf. \cite{Wolibner}, \cite{Holder}, \cite{Kato1967}]
\label{thm:smooth}
Let $\omega_0\in C_c^\infty(\R^2)$ be real with $\supp\omega_0\subset B_{R_0}$.
There is a unique $\omega\in C^\infty(\R\times\R^2)$ such that:
\begin{enumerate}[label=\textup{(\arabic*)}]
\item $\supp\omega(t)\subset B_{R(t)}$ for every $t\in\R$, with
  $R(t)=R_0+(2\|\omega_0\|_{L^\infty}\|\omega_0\|_{L^1}/\pi)^{1/2}|t|$;
\item $v:=\KBS\ast\omega\in C^\infty(\R\times\R^2)$, $v$ and every derivative of
  $v$ are bounded on $[-T,T]\times\R^2$ for every $T>0$, and $\divop v=0$;
\item $\partial_t\omega+v\cdot\nabla\omega=0$ on $\R\times\R^2$ and $\omega(0)=\omega_0$.
\end{enumerate}
Moreover:
\begin{enumerate}[label=\textup{(\arabic*)},resume]
\item the flow $\flow$ of $v$, $\partial_t\flow(t,y)=v(t,\flow(t,y))$,
  $\flow(0,y)=y$, is defined for all $t\in\R$; each $\flow_t=\flow(t,\cdot)$ is
  a $C^\infty$ diffeomorphism of $\R^2$ with $\det D\flow_t\equiv1$; the maps
  $\flow$ and $(t,x)\mapsto\flow_t^{-1}(x)$ lie in $C^\infty(\R\times\R^2)$; and
  $\omega(t)=\omega_0\circ\flow_t^{-1}$; consequently
  $\int F(\omega(t))\dd x=\int F(\omega_0)\dd x$ for every continuous $F$ with
  $F(0)=0$;
\item if $I\ni0$ is an interval and $\tilde\omega\in C^1(I\times\R^2)$ satisfies
  $\tilde\omega(0)=\omega_0$, $\partial_t\tilde\omega+(\KBS\ast\tilde\omega)\cdot\nabla\tilde\omega=0$
  on $I\times\R^2$, and, for each compact subinterval $I'\subset I$,
  $\supp\tilde\omega(t)\subset B_{\tilde R}$ for $t\in I'$ with $\tilde R$
  depending on $I'$, then $\tilde\omega=\omega$ on $I\times\R^2$;
\item if $\omega_0\circ\Orot_\theta=\omega_0$ for a rotation $\Orot_\theta\in SO(2)$,
  then $\omega(t)\circ\Orot_\theta=\omega(t)$ for every $t\in\R$;
\item for every $T>0$ and every $q\in(4/3,\infty)$, the restriction of $\omega$
  to $[0,T)$ lies in $L^\infty_t(L^1\cap L^\infty)_x$ and is a renormalized
  weak solution in $X_q(T)$ with datum $\omega_0$, in the sense of
  Definition~I.\ref{I-def:weak}.
\end{enumerate}
\end{theorem}

The energy estimates of Section~\ref{sec:single}, the construction of
Appendix~\ref{app:smooth} and the proof of Lemma~\ref{lem:Hk-calculus} below use
the following Sobolev inequalities.

\begin{lemma}
\label{lem:embed}
The following embeddings hold on $\R^2$, with explicit constants.
\begin{enumerate}[label=\textup{(E\arabic*)}]
\item For $g\in H^1$, $\|g\|_{L^4}\le\|g\|_{H^1}$.
\item For $g\in H^2$, $g$ has a bounded continuous
  representative, and $\|g\|_{L^\infty}\le\|g\|_{H^2}$.
\item For an integer $k\ge0$ and $g\in H^{2k}$,
  $\|g\|_{H^k}^2\le3^k\|g\|_{L^2}\|g\|_{H^{2k}}$.
\end{enumerate}
\end{lemma}

\begin{proof}
(E1) Let first $g\in C_c^\infty(\R^2)$. Since
$|g(x)|^2=\int_{-\infty}^{x_1}\partial_1|g|^2(t,x_2)\dd t$, we have
\[
|g(x)|^2\le\int_\R|\partial_1|g|^2|(t,x_2)\dd t=:a(x_2),
\]
and likewise
$|g(x)|^2\le\int_\R|\partial_2|g|^2|(x_1,t)\dd t=:b(x_1)$. The product of
these two bounds, integrated over $x$, gives
$\int|g|^4\le\int a\int b$. Since $|\partial_j|g|^2|\le2|g||\partial_jg|$,
the Cauchy-Schwarz inequality, and then the inequality between the arithmetic
and geometric means, give
$\int|g|^4\le4\|g\|_{L^2}^2\allowbreak\|\partial_1g\|_{L^2}\|\partial_2g\|_{L^2}\le2\|g\|_{L^2}^2\|\nabla g\|_{L^2}^2$.
By density of $C_c^\infty$ in $H^1$ the inequality holds on $H^1$. Then, by the
same inequality between the means,
$\|g\|_{L^4}\le2^{1/4}\|g\|_{L^2}^{1/2}\|\nabla g\|_{L^2}^{1/2}\le2^{-1/4}\|g\|_{H^1}$.

(E2) Since $\widehat{\Delta g}=-|\zeta|^2\hat g$, Plancherel's theorem gives
$\|(1+|\zeta|^2)\hat g\|_{L^2}=2\pi\|g-\Delta g\|_{L^2}\le2\sqrt3\,\pi\|g\|_{H^2}$,
where the inequality follows from the triangle inequality for the three terms of
$g-\Delta g$ and the Cauchy-Schwarz inequality.
As $\int_{\R^2}(1+|\zeta|^2)^{-2}\dd\zeta=\pi$, Cauchy-Schwarz gives
$\|\hat g\|_{L^1}\le\sqrt\pi\,\|(1+|\zeta|^2)\hat g\|_{L^2}$. Thus $\hat g\in L^1$, the
Fourier inversion formula gives a bounded continuous representative, and
$\|g\|_{L^\infty}\le(2\pi)^{-2}\|\hat g\|_{L^1}\le\frac{\sqrt3}{2\sqrt\pi}\|g\|_{H^2}\le\|g\|_{H^2}$.

(E3) Put $P_k(\zeta)=\sum_{|\gamma|\le k}\zeta^{2\gamma}$. Every monomial
$\zeta^{2\gamma}$ with $|\gamma|\le k$ occurs in the expansion of
$(1+\zeta_1^2+\zeta_2^2)^k$ exactly once, with a multinomial coefficient between
$1$ and $3^k$; hence $P_k\le(1+|\zeta|^2)^k\le3^kP_k$. By Plancherel's theorem and
Cauchy-Schwarz,
\[
\|g\|_{H^k}^2\le\frac{\int(1+|\zeta|^2)^k|\hat g|^2\dd\zeta}{(2\pi)^2}
\le\frac{\|\hat g\|_{L^2}\|(1+|\zeta|^2)^{k}\hat g\|_{L^2}}{(2\pi)^2}
\le3^k\|g\|_{L^2}\|g\|_{H^{2k}} .\qedhere
\]
\end{proof}

In Section~\ref{sec:single} the perturbation of a ring is supported in a fixed
annulus, and the field that transports it is controlled only on an open
neighborhood of that annulus. The following calculus in $H^k$ is stated for
such data: $w$ is supported in a compact set $K$, and the field $U$ need only be
given on an open set $\mathcal O\supset K$.

\begin{lemma}
\label{lem:Hk-calculus}
Let $k\ge3$ be an integer, $K\subset\R^2$ compact, $\mathcal O\supset K$ open and
bounded, and $\chi_K\in C_c^\infty(\mathcal O)$ with $\chi_K=1$ on a neighborhood of
$K$. There is $C=C(k,\chi_K)$ such that:
\begin{enumerate}[label=\textup{(\arabic*)}]
\item for $w\in H^{k+1}(\R^2)$ with $\supp w\subset K$, a vector field
  $U\in W^{1,\infty}(\mathcal O)\cap H^k(\mathcal O)$, and every multi-index $\alpha$ with
  $|\alpha|\le k$,
  \[
  \|[\partial^\alpha,U\cdot\nabla]w\|_{L^2}\le C\big(\|U\|_{W^{1,\infty}(\mathcal O)}+\|U\|_{H^k(\mathcal O)}\big)\|w\|_{H^k};
  \]
  if moreover $U$ is real and $\divop U=0$ on $\mathcal O$, then
  $|\Re\langle\partial^\alpha(U\cdot\nabla w),\partial^\alpha w\rangle_{L^2}|\le
  C\big(\|U\|_{W^{1,\infty}(\mathcal O)}+\|U\|_{H^k(\mathcal O)}\big)\|w\|_{H^k}^2$, and for
  real $w$ the inner product is real;
\item $\|gh\|_{H^j}\le C\|g\|_{H^j}\|h\|_{H^j}$ for every integer $j$ with
  $2\le j\le k$ and all $g,h\in H^j(\R^2)$;
\item for every integer $j\ge0$, every $\rho>0$ and every $\chi\in C_c^\infty(B_\rho)$
  there is $C'=C'(j,\chi,\rho)$ such that
  $\|\chi\,\KBS\ast w\|_{H^{j+1}}\le C'\|w\|_{H^j}$ for every $w\in H^j(\R^2)$ with
  $\supp w\subset B_\rho$.
\end{enumerate}
Here $[\partial^\alpha,U\cdot\nabla]w=\partial^\alpha(U\cdot\nabla w)-U\cdot\nabla\partial^\alpha w$,
with both terms extended by $0$ outside $\mathcal O$.
\end{lemma}

\begin{proof}
We prove (2) first. Let $g,h\in C_c^\infty(\R^2)$, $|\gamma|\le j$ and
$\beta\le\gamma$. If $\beta=0$ or $\beta=\gamma$, then
$\|\partial^\beta g\,\partial^{\gamma-\beta}h\|_{L^2}\le\|g\|_{H^j}\|h\|_{H^j}$ by
Lemma~\ref{lem:embed}\textup{(E2)}, since $j\ge2$. If
$1\le|\beta|\le|\gamma|-1$, then by H\"older's inequality and
Lemma~\ref{lem:embed}\textup{(E1)},
$\|\partial^\beta g\,\partial^{\gamma-\beta}h\|_{L^2}\le\|\partial^\beta g\|_{H^1}\|\partial^{\gamma-\beta}h\|_{H^1}
\le\|g\|_{H^{|\beta|+1}}\|h\|_{H^{|\gamma|-|\beta|+1}}\le\|g\|_{H^j}\|h\|_{H^j}$.
The Leibniz formula
$\partial^\gamma(gh)=\sum_{\beta\le\gamma}\binom\gamma\beta\partial^\beta g\,\partial^{\gamma-\beta}h$
and $\sum_{\beta\le\gamma}\binom\gamma\beta=2^{|\gamma|}$ give (2) for smooth $g,h$,
with $C=C(k)$. For $g,h\in H^j$ take $g_n,h_n\in C_c^\infty$ converging to
$g,h$ in $H^j$; by the bound just proved, $g_nh_n$ is Cauchy in $H^j$, and by
(E2) it converges to $gh$ in $L^2$, so $gh\in H^j$ and (2) holds.

We now prove (1). Put $\tilde U=\chi_KU$, extended by $0$ outside $\mathcal O$; by the
Leibniz rule, $\|\tilde U\|_{W^{1,\infty}(\R^2)}\le C\|U\|_{W^{1,\infty}(\mathcal O)}$ and
$\|\tilde U\|_{H^k(\R^2)}\le C\|U\|_{H^k(\mathcal O)}$, with $C=C(k,\chi_K)$. Since
$\chi_K=1$ on a neighborhood $V$ of $K$ and $w=0$ outside $K$,
$[\partial^\alpha,U\cdot\nabla]w=[\partial^\alpha,\tilde U\cdot\nabla]w$, and
\[
[\partial^\alpha,\tilde U\cdot\nabla]w=\sum_{0<\beta\le\alpha}\binom\alpha\beta\,
\partial^\beta\tilde U\cdot\nabla\partial^{\alpha-\beta}w :
\]
with $\tilde U_n\in C_c^\infty$ in place of
$\tilde U$, this identity holds in $\mathcal D'$ by the Leibniz formula. Since
$\tilde U\in H^k(\R^2)$ by the bound above, there are such $\tilde U_n$ with
$\tilde U_n\to\tilde U$ in $H^k$. As $w\in H^{k+1}$, we have
$\nabla\partial^{\alpha-\beta}w\in L^2$ for $\beta\le\alpha$; so each product in
the identity, as
well as $\tilde U_n\cdot\nabla w$ and $\tilde U_n\cdot\nabla\partial^\alpha w$,
converges in $L^1$, and the identity passes to the limit. We bound the terms on
the right according to $|\beta|$. Write each of them as
$\sum_l\partial^\beta\tilde U_l\,\partial_l\partial^{\alpha-\beta}w$. If $|\beta|=1$,
its $L^2$ norm is at most $2\|\nabla\tilde U\|_{L^\infty}\|w\|_{H^k}$. If
$|\beta|=|\alpha|\ge2$, it is at most
$\sum_l\|\partial^\alpha\tilde U_l\|_{L^2}\|\partial_lw\|_{L^\infty}\le2\|\tilde U\|_{H^k}\|w\|_{H^3}
\le2\|\tilde U\|_{H^k}\|w\|_{H^k}$, by (E2) and $k\ge3$. If
$2\le|\beta|\le|\alpha|-1$, then by H\"older's inequality and (E1) it is at most
$\sum_l\|\partial^\beta\tilde U_l\|_{H^1}\|\partial_l\partial^{\alpha-\beta}w\|_{H^1}
\le2\|\tilde U\|_{H^{|\beta|+1}}\|w\|_{H^{|\alpha|-|\beta|+2}}\le2\|\tilde U\|_{H^k}\|w\|_{H^k}$,
because $|\beta|+1\le k$ and $|\alpha|-|\beta|+2\le k$. Summing over $\beta$ gives
the commutator bound. For the second claim, put $g=\partial^\alpha w\in H^1$,
which vanishes outside $K$. Then
$\partial^\alpha(U\cdot\nabla w)=[\partial^\alpha,\tilde U\cdot\nabla]w+\tilde U\cdot\nabla g$
and
$\Re\langle\tilde U\cdot\nabla g,g\rangle_{L^2}=\frac12\int\tilde U\cdot\nabla|g|^2\dd x
=-\frac12\int(\divop\tilde U)|g|^2\dd x=0$:
the integration by parts holds because $|g|^2\in W^{1,1}$ and $\tilde U$ is
Lipschitz with compact support, and the last integral vanishes because
$\divop\tilde U=\chi_K\divop U+U\cdot\nabla\chi_K=0$ on $V$ and $g=0$ outside $K$.
Hence the real part is at most
$\|[\partial^\alpha,\tilde U\cdot\nabla]w\|_{L^2}\|g\|_{L^2}$, by the
Cauchy-Schwarz inequality, and the commutator bound gives the claim, because
the $L^2$ norm of $g$ is at most the $H^k$ norm of $w$.

We now prove (3). Let first $w\in C_c^\infty(B_\rho)$ and $u=\KBS\ast w$. By
Lemma~I.\ref{I-lem:BS-basic}(d), and (c) with $R=\rho$,
$\|\nabla u\|_{H^j}=\|w\|_{H^j}$ and $\|u\|_{L^2(B_\rho)}\le2\rho\|w\|_{L^2}$.
For $|\gamma|\le j+1$, in the Leibniz formula
$\partial^\gamma(\chi u)=\sum_{\beta\le\gamma}\binom\gamma\beta\partial^{\gamma-\beta}\chi\,\partial^\beta u$
the term $\beta=0$ has $L^2$ norm at most
$\|\partial^\gamma\chi\|_{L^\infty}\|u\|_{L^2(B_\rho)}$, because $\chi$ vanishes
outside $B_\rho$, and every other term at most
$\|\partial^{\gamma-\beta}\chi\|_{L^\infty}\|\nabla u\|_{H^j}$. This gives (3) for
smooth $w$. For general $w$, $\supp w$ is a compact subset of $B_\rho$, so
mollification gives $w_n\in C_c^\infty(B_\rho)$ with $w_n\to w$ in $H^j$; then
$\chi\KBS\ast w_n$ is Cauchy in $H^{j+1}$ and, by Lemma~I.\ref{I-lem:BS-basic}(c),
converges to $\chi\KBS\ast w$ in $L^2$, so (3) holds for $w$.
\end{proof}
